 We turn now to free groups.  Let $F$ be a free group.  Let $w$ be an element of $F$.  We define the girth of $w$ in $F$ as follows.  Let $g$ be the rank of $F$.  Form a handlebody $H$ of genus $g$.  Identify $F$ with the fundamental group $\pi_{1}(H)$.  Let $\gamma$ be a closed loop in $H$ that gives rise to $w$.  Define the girth of $w$ to be the girth of $\gamma$ in $H$.  We will sometimes write the girth of $w$ as $\girth(w)$.  The following lemma shows this notion is well-defined.  In other words, the girth of an element of a free group does not depend on the choice of identification between $F$ and $\pi_{1}(H)$.  A different identification differs by an automorphism of $\pi_{1}(H)$.

\begin{lemma}
Let $H$ be a handlebody.  Let $\gamma$ be a closed loop in $H$ that represents an element of $\pi_{1}(H)$.  Let $\Phi$ be an automorphism of $\pi_{1}(H)$.  Let $\gamma^{\prime}$ be a closed curve in $H$ that represents the element $\Phi(\gamma)$ in $\pi_{1}(H)$.  The girth of $\gamma$ in $H$ equals the girth of $\gamma^{\prime}$. 
\end{lemma}

\begin{proof}
The proof follows from the fact that automorphisms of a free group can be realized geometrically as homeomorphisms of a handlebody.

Let $F$ be a free group on $g$ letters.  The elementary Nielsen transformations of $F$ generate the group of automorphisms of $F$ (see for instance~\cite[Ch. I.4]{LSch}).  If $F$ is the fundamental group of a handlebody $H$ with generating system $\dd$, the elementary Nielsen transformations can be realized geometrically as a homeomorphism of $H$.  Thus all automorphisms of $F$ can be realized geometrically~\cite{Z}.

Let $n$ be the girth of $\gamma$.  Let $D$ be a girth-realizing disc for $\gamma$ in $H$.  By Lemma~\ref{wedge}, homotope $\gamma$ to intersect $D$ exactly $n$ times.  Realize the automorphism $\Phi$ of $\pi_{1}(H)$ as a homeomorphism of the handlebody $H$.

Apply the homeomorphism $\Phi$ to $H$, which contains $\gamma$ and $D$.  Then $\Phi(\gamma)$ is a curve homotopic to $\gamma^{\prime}$, and $\Phi(\gamma)$ intersects the disc $\Phi(D)$ exactly $n$ times.  Therefore the girth of $\gamma^{\prime}$ is less than or equal to $n$.  The reverse inequality follows by applying the same argument to the curve $\gamma^{\prime}$, in reverse.
\end{proof}

Let $F$ be a free group.  A \emph{generating system} for $F$ is a collection $\dd$ of free generators of $F$.  Let $g$ be the rank of $F$.  Since $F$ is Hopfian, any generating system for $F$ has $g$ elements.  We saw above that a handlebody $H$ with disc system $\dd$ gives rise to a generating system for $\pi_{1}(H)$.  

Let $\dd$ be a generating system for $F$.  Let $w$ be an element of $F$.  To parallel the definition above, we define $w\hi\dd$ as follows.  Let $H$ be a handlebody with disc system $\dd$.  Then $\pi_{1}(H)$ is naturally identified with $F$.  Choose a closed curve $\gamma$ that gives rise to $w$.  Define $w\hi\dd$ to equal $\gamma\hi\dd$.  In Section~\ref{computecurve}, we call $w\hi\dd$ the complexity of $\dd$ with respect to $w$.  The following lemma gives another meaning to $w\hi\dd$.

\begin{lemma}
\label{cyclic}
Let $F$ be a free group.  Let $\dd$ be a generating system for $F$.  Let $w$ be any element of $F$.  Then $w\hi\dd$ equals the length of $w$ written as a cyclically reduced word in the generators $\dd$ and their inverses.
\end{lemma}

Note that any element of the free group on $g$ letters has a unique cyclically reduced form, up to cyclic permutation (see~\cite[Sec. 1.4]{MKS}).  However, this is not needed for the proof.

\begin{proof}[Proof of Lemma~\ref{cyclic}]
Represent $F$ as the fundamental group of a handlebody $H$ with oriented disc system $\dd$.  The fundamental group of $H$ is in natural correspondence with $F$.  Consider any closed curve $\gamma$ in $H$.  The pattern of intersections with $\dd$ gives rise to a cyclic word $w$ in the generators $\dd$.  By Lemma~\ref{wedge}, the quantity $w\hi\dd$ equals the minimum intersection number of $\dd$ and a curve homotopic to $\gamma$.

It suffices to show that if the length of the cyclic word $w$ can be reduced, then $|\gamma\cap\dd|$ can be reduced.  Assume $w$ can be cyclically reduced in length.  By~\cite[Sec. 1.4]{MKS}, there is an adjacent pair of letters $ab$ in $w$ where $a=b^{-1}$.  Choose such a pair $ab$.  Consider the tree $\Gamma$ associated to $\dd$.  Lift $\gamma$ to $\Gamma$.  The pair $ab$ corresponds to a backtracking in the lift $\gamma$.  As in the proof of Lemma~\ref{wedge}, this backtracking gives rise to an elementary homotopy of $\gamma$ that decreases the intersection with $\dd$.  See Figure~\ref{homotopy 1} for the picture.
\end{proof}

Let $w$ be an element of the free group $F$.   By the following lemma, girth satisfies the group multiplicative property $\girth(w^{m})=m \girth(w)$.

\begin{lemma}
\label{power}
For $w$ in a free group $F$, $\girth(w^m)=m \girth(w)$.
\end{lemma}

This would follow from Lemma~\ref{cyclic} if girth-realizing discs were always non-separating.  We introduce the concept of a complete disc system for the proof.  A \emph{complete disc system} is a disjoint collection of essential discs in $H$ that cut $H$ into solid pairs of pants.  If $g$ is the genus of $H$, any complete disc system has $3g-3$ discs and cuts $H$ into $2g-2$ pairs of pants. 

\begin{proof}
Assume $w$ is nontrivial.  Let $g$ be the rank of $F$.  Form a handlebody $H$ of genus $g$.  Identify $F$ with the fundamental group $\pi_{1}(H)$.  Let $\gamma$ be a closed loop in $H$ that represents the element of $\pi_{1}(H)$ corresponding to $w$.  Then the girth of $w$ equals the girth of $\gamma$, and the girth of $w^{m}$ equals the girth of $\gamma^{m}$.

Let $n_{1}$ be the girth of $\gamma$.  Let $D_{1}$ be a girth-realizing disc for $\gamma$.  Freely homotope $\gamma$ to intersect $D_{1}$ exactly $n_{1}$ times.  Then $m$ copies of $\gamma$ intersect $D_{1}$ exactly $mn_{1}$ times.  Thus, $\girth(\gamma^m)\leq mn_{1}=m \girth(\gamma)$.

We now prove the reverse inequality.  Let $n_{2}$ be the girth of $\gamma^{m}$.  Let $D_{2}$ be a girth-realizing disc for $\gamma^m$.  Then $\gamma^{m}\hi D_{2}$ equals $n_{2}$.  Choose a complete disc system $\dd$ that contains $D_{2}$.  Freely homotope $\gamma^m$ to minimize intersection with $\dd$.  Since $\dd$ is a complete disc system and $\gamma^{m}$ is nontrivial, we have $|\gamma^{m}\cap\dd|>0$.  By Lemma~\ref{subcollection}, $\gamma^{m}$ also minimizes geometric intersection with $D_{2}$.  Furthermore, by Lemma~\ref{wedge}, the curve $\gamma^{m}$ intersects $D_{2}$ exactly $n_{2}$ times.

Let $\Gamma$ be the tree associated to $\dd$.  The graph $\Gamma$ is a trivalent graph.  By Lemma~\ref{geodesic}, the curve $\gamma^{m}$ lifts to a path $P$ in $\Gamma$ without backtrackings.  Since $\gamma^{m}$ is not disjoint from $\dd$, the path $P$ is an infinite geodesic path.  The action of $F$ on $\hh$ via deck transformations induces an action of $F$ on $\Gamma$.  The element $w^{m}$ leaves $P$ invariant.  Indeed, $P/\langle w^{m}\rangle$ is a path in $\Gamma/\langle w^{m}\rangle$ that crosses the edge corresponding to $D_{2}$ exactly $n_{2}$ times.

To complete the proof, it suffices to show that $w$ also leaves $P$ invariant.  For if this were the case, then $P/\langle w\rangle$ would be a circuit crossing the edge corresponding to $D_{2}$ exactly $n_{2}/m$ times.  The sequence of edges crossed by this circuit would define a loop in $H$ that is freely homotopic to $\gamma$ and that crosses $D_{2}$ exactly $n_{2}/m$ times.  Hence, $\girth(w)\leq n_{2}/m=\girth(\gamma^{m})/m$.

To see that $w$ leaves $P$ invariant, observe that $w^{m}$ leaves the path $w(P)$ invariant.  By a theorem of Tits~\cite[Prop. 24]{Se}, $w(P)$ must contain the infinite geodesic $P$.  Thus $w(P)$ and $P$ are actually the same.
\end{proof}